% !TEX root = Lecture_Koszul.tex
\section{The quadratic dual}\label{sec:quadratic-dual}


From now on, we require the Assumption~\ref{ass:basic-semisimple}: $A_0=S=\bigoplus_{i=1}^n\Bbbk e_i$.  The general case when this assumption is removed will be dealt with in the advanced sections.


\subsection{Orthogonal complements and quadratic duality}

\begin{proposition}[Duals of tensor products]\label{prop:D-tensor-duals}
For finite-dimensional bimodules ${}_SV_S$ and ${}_SW_S$, there is a natural bimodule isomorphism
\begin{equation}\label{eq:D-tensor-duals}
\theta:D(W)\otimes_S D(V)\xrightarrow{\sim}D(V\otimes_S W),\qquad
\theta(\lambda\otimes\mu)(v\otimes w)=\sum_{i=1}^n\mu(ve_i)\,\lambda(e_iw).
\end{equation}
\comm{add bi-naturality}
\end{proposition}

\begin{supplementary}
\begin{proof}
By \eqref{eq:linear-bimodule-dual-action}, $(\lambda e_k)(w)=\lambda(e_kw)$ and $(e_k\mu)(v)=\mu(ve_k)$.
Hence, for every $k$, replacing $v\otimes e_kw$ by $ve_k\otimes w$, or $\lambda\otimes e_k\mu$ by $\lambda e_k\otimes\mu$, leaves the right-hand side equal to $\mu(ve_k)\lambda(e_kw)$.
By linearity the formula respects both tensor relations, so $\theta$ is well defined.
For $s,t\in S$, the same rule gives $\theta(s\lambda\otimes\mu t)(v\otimes w)=\theta(\lambda\otimes\mu)(tv\otimes ws)$, so $\theta$ is a bimodule map; it commutes with precomposition by bimodule maps, so it is natural.
For bijectivity, Lemma~\ref{lem:idempotent-tensors} and \eqref{eq:dual-blocks} decompose both sides into summands indexed by $(i,j,\ell)$, on which $\theta$ is
\[
D(e_jWe_\ell)\otimes_\Bbbk D(e_iVe_j)\longrightarrow D(e_iVe_j\otimes_\Bbbk e_jWe_\ell),\qquad
\lambda\otimes\mu\longmapsto\bigl(v\otimes w\mapsto\mu(v)\lambda(w)\bigr).
\]
This map sends tensor products of dual-basis vectors to the dual basis of the tensor-product basis, so it is an isomorphism.
\end{proof}
\end{supplementary}
\begin{remark}
Note that the use of $e_i$ is necessary, we cannot just take $\theta(\lambda\otimes \mu)(v\otimes w)=\mu(v)\lambda(w)$ (take for example $v\otimes w$ as $e_1\otimes e_2$ to see this).
\end{remark}

\begin{advanced}
\begin{remark}[General $S$]
For an arbitrary finite-dimensional semisimple $S$, there are no coordinate idempotents to sum over.
After choosing a symmetrising form $\tau$ and bases $(x_k)$, $(y_k)$ of $S$ with $\tau(x_ky_l)=\delta_{kl}$, the sum $\sum_k\mu(vx_k)\lambda(y_kw)$ replaces \eqref{eq:D-tensor-duals}.
Alternatively, the one-sided $S$-duals give the choice-free isomorphisms $\theta_r,\theta_\ell$ in \eqref{eq:tensor-duals} (advanced).
\end{remark}
\end{advanced}

\begin{definition}\label{def:D-orthogonal}
For a sub-bimodule $U$ of a finite-dimensional bimodule ${}_SM_S$, its \defn{orthogonal complement} is
\[
U^\perp:=\{\lambda\in D(M)\mid\lambda(u)=0\text{ for all }u\in U\}.
\]
\end{definition}

It is a sub-bimodule of $D(M)$: for $s,t\in S$ and $u\in U$, we have $(s\lambda t)(u)=\lambda(tus)=0$ since $tus\in U$.
In fact, by definition we have.
\[U^\perp = \Ker(D\iota: DM\twoheadrightarrow DU)) \cong D(M/U)\]
for $\iota:U\hookrightarrow M$ the canonical inclusion of $U$ as sub-bimodule of $M$.  In other words, applying $D$ to the canonical short exact sequence 
\[
0\to U\to M\to M/U \to 0 \quad \text{yields} \quad 0\to U^\perp \to DM\to DU\to 0
\]
in $\mod S^\en$.


\begin{advanced}
For an arbitrary finite-dimensional semisimple $S$, the one-sided $S$-duals give the following choice-free orthogonal complements.

\begin{definition}\label{def:orthogonals}
For an $S$-sub-bimodule $U$ of a finite-dimensional $S$-bimodule $M$, its \defn{right orthogonal complement} and \defn{left orthogonal complement} are
\[
\begin{aligned}
U^\perp&:=\{f\in M^*\mid f(u)=0\text{ for all }u\in U\},\\
{}^\perp U&:=\{g\in{}^*M\mid g(u)=0\text{ for all }u\in U\}.
\end{aligned}
\]
\end{definition}

Both orthogonal complements are $S$-sub-bimodules.
Indeed, for $s,t\in S$ and $u\in U$, the action on $S$-dual \eqref{eq:bimodule-dual-actions} reads
\[
(sft)(u)=sf(tu)=0\quad(f\in U^\perp),\qquad
(sgt)(u)=g(us)t=0\quad(g\in{}^\perp U),
\]
since $tu,us\in U$.
The orthogonal complements are defined without choosing a symmetrising form, but it is possible to realise, through the natural isomorphism associated to a chosen symmetrising form, the orthogonal complements as a sub-bimodule of $D(M)$.

If $S$ is equipped with a symmetrising form $\tau:S\to\Bbbk$, the comparison maps $\alpha_M,\beta_M$ of \eqref{eq:dual-comparison} satisfy
\begin{equation}\label{eq:orthogonal-comparison}
\alpha_M(U^\perp)=\{\lambda\in D(M)\mid\lambda(u)=0\text{ for all }u\in U\}=\beta_M({}^\perp U)
.
\end{equation}
The middle space is the orthogonal complement of Definition~\ref{def:D-orthogonal}; it contains $\alpha_M(U^\perp)$ and $\beta_M({}^\perp U)$: if $f\in U^\perp$ and $g\in{}^\perp U$, then $\tau\circ f$ and $\tau\circ g$ vanish on $U$.
Conversely, $\alpha_M(f)(U)=0$ gives
\[
\tau(f(u)s)=\tau(f(us))=0
\quad\text{for all }s\in S,\ u\in U.
\]
Non-degeneracy of $\tau$ forces $f(u)=0$ for every $u\in U$.
For the left dual, $\tau(sg(u))=\tau(g(su))=0$ gives $g(u)=0$ instead.
Bijectivity of $\alpha_M,\beta_M$ in Theorem~\ref{prop:dual-comparison} proves \eqref{eq:orthogonal-comparison}.
\end{advanced}

\begin{example}\label{ex:symmetric-orthogonals}
We compute the orthogonal complements of the relation spaces of $\SymAlg(V)$ and $\bigwedge W$ for $\dim_{\Bbbk}V=2=\dim_{\Bbbk}W$.
Take $V=\Bbbk x_1\oplus\Bbbk x_2$ and $W=\Bbbk y_1\oplus\Bbbk y_2$ in Example~\ref{ex:polynomial-exterior}.
The relation spaces in $V\otimes_\Bbbk V$ and $W\otimes_\Bbbk W$, respectively, are
\[
\begin{aligned}
R_{\SymAlg(V)}&=\Bbbk(x_1\otimes x_2-x_2\otimes x_1),\\
R_{\bigwedge W}&=\Bbbk(y_1\otimes y_1)\oplus\Bbbk(y_2\otimes y_2)\oplus\Bbbk(y_1\otimes y_2+y_2\otimes y_1).
\end{aligned}
\]
Let $x_1^*,x_2^*$ and $y_1^*,y_2^*$ be the dual bases of $D(V)$ and $D(W)$, respectively.
By Proposition~\ref{prop:D-tensor-duals} with $S=\Bbbk$, where $\theta(\lambda\otimes\mu)(v\otimes w)=\mu(v)\lambda(w)$,
\[
\begin{aligned}
R_{\SymAlg(V)}^\perp&=\operatorname{span}_\Bbbk\{x_1^*\otimes x_1^*,\,x_2^*\otimes x_2^*,\,x_1^*\otimes x_2^*+x_2^*\otimes x_1^*\},\\
R_{\bigwedge W}^\perp&=\Bbbk(y_1^*\otimes y_2^*-y_2^*\otimes y_1^*).
\end{aligned}
\]
The first space consists of the exterior relations on $D(V)$; the second consists of the symmetric relations on $D(W)$.
\end{example}


\subsection{Generators and orthogonal relations}

Fix a quadratic presentation $A=T_S(V)/(R)$ as in Definition~\ref{def:quadratic}.

%\begin{equation}\label{eq:quadratic-pairings}
%\begin{aligned}
%\theta_r(g\otimes f)(v\otimes w)&=g(f(v)w),
%& f,g&\in V^*,\\
%\theta_\ell(g\otimes f)(v\otimes w)&=f(vg(w)),
%& f,g&\in{}^*V.
%\end{aligned}
%\end{equation}
%Here $v,w\in V$. Both isomorphisms reverse the order of the factors.

By Proposition~\ref{prop:D-tensor-duals}, there is a bimodule isomorphism $\theta:D(V)\otimes_S D(V)\xrightarrow{\sim}D(V\otimes_S V)$, given by
\[
\theta(\lambda\otimes\mu)(v\otimes w)=\sum_{i=1}^n\mu(ve_i)\,\lambda(e_iw)
\qquad(\lambda,\mu\in D(V),\ v,w\in V).
\]
We use $\theta$ to identify $D(V)\otimes_S D(V)$ with $D(V\otimes_S V)$.
In particular, the orthogonal complement $R^\perp\subset D(V\otimes_S V)$ of Definition~\ref{def:D-orthogonal} becomes a sub-bimodule of $D(V)\otimes_S D(V)$.

\begin{definition}\label{def:quadratic-duals}
The \defn{quadratic dual} of $A$ is
\begin{equation}\label{eq:quadratic-duals}
A^!:=T_S(D(V))/(R^\perp),
\qquad
R^\perp=\{F\in D(V)\otimes_S D(V)\mid\theta(F)(R)=0\}.
\end{equation}
\end{definition}

Note that $D(V)$ is placed in degree $1$ in \eqref{eq:quadratic-duals} -- in contrast to the graded dual $\bbD(V)$, which is placed in degree $-1$.
By \eqref{eq:dual-blocks}, dualising reverses the idempotents of each generator: $e_jD(V)e_i\cong D(e_iVe_j)$.

Note also that the quadratic dual is independent to the choice of quadratic presentation of $A$, up to $A_0$-preserving isomorphisms of graded algebras. 
Indeed, by Proposition~\ref{prop:quadratic-basic}, a quadratic presentation yields 
\[S\cong A_0, V\cong A_1, \text{ and }R\cong \Ker(\mu_2:A_1\otimes_{A_0}A_1 \to A_2).
\]
Through these isomorphisms, functoriality of $D$ and Proposition~\ref{prop:D-tensor-duals} give
\[
D(V)\cong D(A_1),\qquad
R^\perp\cong(\Ker\mu_2)^\perp\subset D(A_1)\otimes_{A_0}D(A_1).
\]
By the universal property of tensor algebra (Proposition~\ref{prop:tensor-universal}), we can extend $D(V)\cong D(A_1)$ to a tensor algebra isomorphism, which descends to the quotients by $(R^\perp)$ and $((\Ker\mu_2)^\perp)$.

\begin{proposition}\label{prop:dual-low-degrees}
The algebra $A^!$ is quadratic and
\[
A^!_0=S,\qquad A^!_1=D(V),\qquad A^!_2\cong D(R).
\]
In particular, $\dim_\Bbbk A^!_2=\dim_\Bbbk R$.
\end{proposition}

\begin{proof}
To see that $A^!$ is quadratic, apply $D$ to the canonical short exact sequence $0\to R\to V\otimes_S V\to A_2\to 0$ of $S$-bimodules, which yields by exactness of $D$ a short exact sequence
\begin{equation}\label{eq:restriction-relations}
0\longrightarrow R^\perp\longrightarrow D(V)\otimes_S D(V)
\xrightarrow{F\mapsto\theta(F)|_R}D(R)\longrightarrow0,
\end{equation}
of $S$-bimodule.
In particular, we get $A_2^!\cong D(R)$ by Proposition \ref{prop:quadratic-basic}, and the remaining claim follows.
\end{proof}
%
%Adding relations has the opposite effect after dualising. For $S$-sub-bimodules $R\subset R'\subset V\otimes_S V$,
%\[
%R\subset R'\quad\Longrightarrow\quad(R')^\perp\subset R^\perp
%\quad\Longrightarrow\quad
%\bigl(T_S(V)/(R')\bigr)^!\onto\bigl(T_S(V)/(R)\bigr)^!.
%\]
%The displayed surjection fixes $S$ and is induced by $\id_{V^*}$ on the generators.

\subsection{Recovering the original algebra}

% ChangeProposal9 consistency update: forward bidual maps now use ev
% (evaluation), as in lem:S-bidual. The left/right dual conventions and
% mathematical arguments are unchanged.

\begin{lemma}\label{lem:double-orthogonal}
Let $U$ be a sub-bimodule of a finite-dimensional bimodule ${}_SM_S$, and identify $M$ with $D(D(M))$ by evaluation (Lemma~\ref{lem:linear-bidual}).
Then $(U^\perp)^\perp=U$.
\end{lemma}
\begin{proof}
This follows from 
\[(U^\perp)^\perp \cong D\Big(\frac{DM}{U^\perp}\Big)\cong D\Big(\frac{DM}{D(M/U)}\Big)\cong D D (U)\cong U.
\]
\end{proof}
\begin{supplementary}
\begin{proof}[By direct calculatin]
If $u\in U$ and $\lambda\in U^\perp$, then $\mathrm{ev}_M(u)(\lambda)=\lambda(u)=0$; hence $U\subset(U^\perp)^\perp$.
If $m\notin U$, extend a basis of $U$ by $m$ to a basis of $M$; the coordinate functional of $m$ lies in $U^\perp$ and does not vanish at $m$, so $m\notin(U^\perp)^\perp$.
\end{proof}
\end{supplementary}

\begin{theorem}\label{thm:quadratic-involution}
Quadratic duality is an involution: evaluation induces an $S$-preserving isomorphism of graded algebras
\[
(A^!)^!\cong A.
\]
\end{theorem}

\begin{proof}
\emph{Generators.} Evaluation $\mathrm{ev}_V:V\xrightarrow{\sim}D(D(V))$ is a bimodule isomorphism by Lemma~\ref{lem:linear-bidual} and \eqref{eq:linear-bimodule-dual-action}; it identifies the generators of $A$ and $(A^!)^!$.

Now, universal property of tensor algebra (Proposition \ref{prop:tensor-universal}) implies that $\mathrm{ev}$ induces a graded algebra homomorphism $A\to (A^!)^!$.

\emph{Relations.} 
It remains to show that $\mathrm{ev}$ induces an isomorphism $(R^\perp)^perp$ with $R$, but this follows from Lemma \ref{lem:double-orthogonal}.
\old{
The relations of $(A^!)^!$ are $(R^\perp)^\perp\subset D(D(V))\otimes_S D(D(V))$, formed with the map $\theta$ for the bimodule $D(V)$.
By \eqref{eq:linear-bimodule-dual-action}, $(\lambda e_i)(w)=\lambda(e_iw)$ and $(e_i\mu)(v)=\mu(ve_i)$, so for $v,w\in V$ and $\lambda,\mu\in D(V)$,
\begin{equation}\label{eq:bidual-tensor-square}
\theta\bigl(\mathrm{ev}_V(v)\otimes\mathrm{ev}_V(w)\bigr)(\lambda\otimes\mu)
=\sum_{i=1}^n\mathrm{ev}_V(w)(\lambda e_i)\,\mathrm{ev}_V(v)(e_i\mu)
=\theta(\lambda\otimes\mu)(v\otimes w).
\end{equation}
Thus $\mathrm{ev}_V\otimes\mathrm{ev}_V$ is the evaluation identification of $V\otimes_S V$ with the dual of $D(V\otimes_S V)\cong D(V)\otimes_S D(V)$.
Lemma~\ref{lem:double-orthogonal}, applied to $R\subset V\otimes_S V$, identifies $(R^\perp)^\perp$ with $R$.
The presentations \eqref{eq:quadratic-duals} therefore give an $S$-preserving isomorphism of graded algebras $(A^!)^!\cong A$.}
\end{proof}

\begin{advanced}
\begin{remark}[General $S$]
For an arbitrary finite-dimensional semisimple $S$, the map $\theta$, and hence Definition~\ref{def:quadratic-duals}, needs a choice of symmetrising form.
The choice-free version uses the right and left $S$-duals: the right and left quadratic duals below satisfy ${}^!(A^!)\cong A\cong({}^!A)^!$ with no choice (Theorem~\ref{thm:quadratic-bidual}, advanced).
\end{remark}

\emph{Quadratic duals via one-sided $S$-duals.}
For an arbitrary finite-dimensional semisimple $S$, we use the isomorphisms $\theta_r,\theta_\ell$ of \eqref{eq:tensor-duals} instead of $\theta$.

\begin{definition}\label{def:S-quadratic-duals}
The \defn{right quadratic dual} and \defn{left quadratic dual} of $A$ are
\begin{equation}\label{eq:S-quadratic-duals}
\begin{aligned}
A^!&:=T_S(V^*)/(R^\perp),
&R^\perp&=\{F\in V^*\otimes_S V^*\mid\theta_r(F)(R)=0\},\\
{}^!A&:=T_S({}^*V)/({}^\perp R),
&{}^\perp R&=\{F\in{}^*V\otimes_S{}^*V\mid\theta_\ell(F)(R)=0\}.
\end{aligned}
\end{equation}
\end{definition}

Under Assumption~\ref{ass:basic-semisimple}, the map $\alpha_V:V^*\to D(V)$, $f\mapsto\tau_S\circ f$, of Theorem~\ref{prop:dual-comparison} satisfies $\theta\circ(\alpha_V\otimes\alpha_V)=\alpha_{V\otimes_SV}\circ\theta_r$: both sides send $g\otimes f$ to $v\otimes w\mapsto\tau_S(g(f(v)w))$, because $\sum_i\tau_S(se_i)e_i=s$ for $s\in S$.
Hence $\alpha_V$ induces an isomorphism from the right quadratic dual onto the quadratic dual of Definition~\ref{def:quadratic-duals}; this is why we use the same symbol $A^!$.

Our notation is swapped relative to \cite[\S2.7 and Definition~2.8.1]{BGS96}: their $V^*$ uses left $S$-linear maps and their ${}^*V$ uses right $S$-linear maps.
Thus our $A^!$ is their ${}^!A$, and our ${}^!A$ is their $A^!$.
The names agree: they also call ${}^!A$ the right dual and $A^!$ the left dual.

\begin{lemma}\label{lem:S-double-orthogonal}
Let $U\subset M$ be an $S$-sub-bimodule of a finite-dimensional bimodule $ {}_SM_S$.
Under $M\cong{}^*(M^*)$ and $M\cong({}^*M)^*$ from the right and left $S$-dualities in Lemma~\ref{lem:S-bidual},
\[
{}^\perp(U^\perp)=U,\qquad ({}^\perp U)^\perp=U.
\]
\end{lemma}

\begin{proof}
For $u\in U$ and $f\in U^\perp$, evaluation $\mathrm{ev}_M:M\to{}^*(M^*)$ gives $\mathrm{ev}_M(u)(f)=f(u)=0$, so $U\subset{}^\perp(U^\perp)$ under Lemma~\ref{lem:S-bidual}.

Conversely, suppose that there is some $m\notin U$, semisimplicity of $S$ implies (by Corollary~\ref{cor:semisimple-summands}) that $M/U\hookrightarrow S^{\oplus n}$ as a right $S$-module for some $n$.
Consider the composition $M\onto M/U\hookrightarrow S^{\oplus n}\twoheadrightarrow S$ which sends $m$ to a non-zero element in a direct summand of $S^{\oplus n}$.
Then this is a right $S$-linear map in $\Hom_S(M,S)=M^*$ satisfying $f(U)=0$ and $f(m)\neq0$.
Thus $f\in U^\perp$ and $m\notin{}^\perp(U^\perp)$.
The proof of $({}^\perp U)^\perp=U$ is similar.
\end{proof}

\begin{theorem}\label{thm:quadratic-bidual}
Left and right quadratic duality are canonically inverse:
\[
{}^!(A^!)\cong A\cong({}^!A)^!
\quad\text{by }S\text{-preserving isomorphisms of graded algebras}.
\]
\end{theorem}

\begin{proof}
\emph{Generators.} The right $S$-duality in Lemma~\ref{lem:S-bidual} gives
$\mathrm{ev}_V:V\xrightarrow{\sim}{}^*(V^*)$, $v\mapsto(f\mapsto f(v))$, identifying the generators of $A$ and $ {}^!(A^!)$.
It preserves both $S$-actions by \eqref{eq:bimodule-dual-actions}.

\emph{Relations.} The right $S$-action on $V^*$ is $(gs)(w)=g(sw)$ for all $g\in V^*$, $s\in S$, and $w\in V$.
Applying \eqref{eq:tensor-evaluations} first to $V^*\otimes_S V^*$ and then to $V\otimes_S V$ gives, for all $v,w\in V$ and $f,g\in V^*$,
\begin{equation}\label{eq:S-bidual-tensor-square}
\begin{aligned}
\theta_\ell(\mathrm{ev}_V(v)\otimes\mathrm{ev}_V(w))(g\otimes f)
&=\mathrm{ev}_V(w)(g f(v))=g(f(v)w)\\
&=\mathrm{ev}_{V\otimes_S V}(v\otimes w)(\theta_r(g\otimes f)).
\end{aligned}
\end{equation}
Thus $\mathrm{ev}_V\otimes\mathrm{ev}_V$ is the bidual identification for $V\otimes_S V$ under anti-distributivity of the $S$-dual \eqref{eq:tensor-duals}.
Lemma~\ref{lem:S-double-orthogonal}, applied to $R\subset V\otimes_S V$, identifies the relations $ {}^\perp(R^\perp)$ of $ {}^!(A^!)$ with $R$.
The presentations \eqref{eq:S-quadratic-duals} therefore give an $S$-preserving isomorphism of graded algebras $ {}^!(A^!)\cong A$.
Interchanging left and right gives $({}^!A)^!\cong A$.
\end{proof}

\begin{proposition}\label{prop:quadratic-dual-sides}
A choice of symmetrising form $\tau:S\to\Bbbk$ gives $S$-preserving isomorphisms of graded algebras $A^!\cong{}^!A$ and $(A^!)^!\cong A$.
\end{proposition}

\begin{proof}
The maps $\alpha_V(f)=\tau\circ f$ and $\beta_V(g)=\tau\circ g$ in \eqref{eq:dual-comparison} give a bimodule isomorphism $\beta_V^{-1}\alpha_V:V^*\to{}^*V$.
For $f,g\in V^*$, write $f'=\beta_V^{-1}\alpha_V(f)$ and $g'=\beta_V^{-1}\alpha_V(g)$. Then
\[
\begin{aligned}
\tau\bigl(f'(vg'(w))\bigr)
&=\tau\bigl(f(v)g'(w)\bigr)
=\tau\bigl(g'(f(v)w)\bigr)\\
&=\tau\bigl(g(f(v)w)\bigr)
\qquad(v,w\in V).
\end{aligned}
\]
By \eqref{eq:orthogonal-comparison}, $(\beta_V^{-1}\alpha_V)^{\otimes_S2}$ therefore identifies $R^\perp$ with ${}^\perp R$.
The bimodule map $\beta_V^{-1}\circ\alpha_V$ extends by Proposition~\ref{prop:tensor-universal} to an isomorphism $T_S(V^*)\to T_S({}^*V)$ and hence descends to $A^!\cong{}^!A$.
Apply the same comparison to $A^!$ and then Theorem~\ref{thm:quadratic-bidual} to obtain $(A^!)^!\cong{}^!(A^!)\cong A$.
\end{proof}

For an arbitrary finite-dimensional semisimple $S$, the isomorphism $A^!\cong{}^!A$ of Proposition~\ref{prop:quadratic-dual-sides} depends on $\tau$; changing $\tau$ changes it by the central unit of Remark~\ref{rem:general-S-comparison}.
Theorem~\ref{thm:quadratic-bidual} involves no choice.
\end{advanced}

\subsection{Tensor algebras and monomial relations}

\begin{advanced}
\begin{example}\label{ex:tensor-square-zero-dual}
Quadratic duality exchanges having no relations with setting every product of two generators equal to zero.
That is, taking $R=0$ or $R=V\otimes_S V$ in \eqref{eq:quadratic-duals} gives
\[
\begin{aligned}
\bigl(T_S(V)\bigr)^!
&=T_S(D(V))/(D(V)\otimes_S D(V))\cong S\oplus D(V),\\
\bigl(T_S(V)/(V\otimes_S V)\bigr)^!&=T_S(D(V)).
\end{aligned}
\]
In $S\oplus D(V)$, $D(V)^2=0$ and $(s,f)(t,g)=(st,sg+ft)$; this construction is also known as \emph{trivial extension of $S$ by $D(V)$}.
\end{example}

For the rest of this subsection, take semisimple $S=\bigoplus_{i=1}^n\Bbbk e_i$ with coordinate idempotents $e_i$.
Choose a basis $\mathcal B_{ij}$ of each $e_iVe_j$ and let $\mathcal B=\bigcup_{i,j}\mathcal B_{ij}$.
A non-empty word $b_1\cdots b_d$ is \defn{composable} if $b_t\in\mathcal B_{i_{t-1},i_t}$ for $1\leq t\leq d$ and some $i_0,\ldots,i_d\in\{1,\ldots,n\}$.
(Thus adjacent factors have matching local units.)
By Lemma~\ref{lem:idempotent-tensors}, the composable words form a basis of the positive degree part of $T_S(V)$.

\begin{definition}\label{def:monomial}
A homogeneous two-sided ideal $I\subset T_S(V)$ with $I_0=I_1=0$ is \defn{monomial} with respect to $\mathcal B$ if it is generated by a set $\mathcal R$ of composable words in $\mathcal B$.
Each word in $\mathcal R$ then has length at least $2$, and $T_S(V)/I$ is called a \defn{monomial algebra}.
If every word in $\mathcal R$ has length $2$, then $T_S(V)/(R)$ with $R=\operatorname{span}_\Bbbk\mathcal R$ is a \defn{monomial quadratic presentation}.
\end{definition}

For example, $\Bbbk A_4/(e_{14})$ in Example~\ref{ex:cubic-quotient} is a monomial algebra, with $\mathcal R=\{e_{12}e_{23}e_{34}\}$, but it is not quadratic.

For $b\in\mathcal B_{ij}$, define the functional $b^*\in D(V)$ by its values on $\mathcal B$:
\begin{equation}\label{eq:dual-basis-idempotents}
b^*(c)=\delta_{bc}\quad(c\in\mathcal B),
\qquad e_jb^*e_i=b^*,
\end{equation}
where the second identity follows from \eqref{eq:dual-blocks}.
The elements $b^*$ form the dual $\Bbbk$-basis of $D(V)$; the left and right local units of each basis element are reversed.

\begin{proposition}\label{prop:monomial-dual}
Let $A=T_S(V)/I$ be a monomial algebra, with $I$ generated by $\mathcal R$.
\begin{enumerate}
\item A basis of $A$ consists of the $e_i$ and the non-empty composable words containing no element of $\mathcal R$ as a consecutive subword.
\item For a monomial quadratic presentation, the quadratic dual is monomial, with
\begin{equation}\label{eq:monomial-dual}
R^\perp=\operatorname{span}_\Bbbk
\{c^*\otimes_S b^*\mid bc\text{ is composable and }bc\notin\mathcal R\}.
\end{equation}
\end{enumerate}
\end{proposition}

\begin{proof}
The ideal $I$ is spanned by the products $urw$ with $r\in\mathcal R$ and $u,w$ each a composable word or an idempotent $e_i$.
The non-zero products are exactly the composable words containing an element of $\mathcal R$ as a consecutive subword.
Removing those basis words proves (1).
For $b\in\mathcal B_{ij}$ and $c\in\mathcal B_{j\ell}$, \eqref{eq:D-tensor-duals} and \eqref{eq:dual-basis-idempotents} give
\[
\theta(c^*\otimes b^*)(b\otimes c)=\sum_{k=1}^n b^*(be_k)\,c^*(e_kc)=b^*(b)\,c^*(c)=1.
\]
The functional vanishes on every other basis tensor.
By Lemma~\ref{lem:idempotent-tensors}, the composable $c^*\otimes b^*$ form a basis of $D(V)\otimes_S D(V)$.
Evaluating a linear combination on each $b\otimes c\in\mathcal R$ forces precisely the coefficient of $c^*\otimes b^*$ to be zero.
Thus the annihilator is the span of the complementary basis tensors, proving (2).
\end{proof}

\begin{example}\label{ex:An-quadratic-dual}
We apply Example~\ref{ex:tensor-square-zero-dual} to the tensor algebra description of $\Bbbk A_n$ in Example~\ref{ex:tensor-An}.
Take $A=\Bbbk A_n$, $S=\bigoplus_{i=1}^n\Bbbk e_{ii}$, and $V=\bigoplus_{i=1}^{n-1}\Bbbk e_{i,i+1}$.
There are no quadratic relations, so $A^!$ has the reversed generators $e_{i,i+1}^*\in e_{i+1,i+1}D(V)e_{ii}$ and every composable product of two generators is zero.
With $C$ as in Example~\ref{ex:hilbert-triangular}, this gives
\[
h_{A^!}(t)=I_n+tC^{\mathsf T}.
\]
Here ${}^{\mathsf T}$ denotes transpose: $(C^{\mathsf T})_{ij}=C_{ji}$.
For $n=3$,
\[
(\Bbbk A_3)^!\cong\bigl(\Bbbk A_3/(e_{13})\bigr)^\op,
\qquad
\bigl(\Bbbk A_3/(e_{13})\bigr)^!\cong(\Bbbk A_3)^\op.
\]
These are graded isomorphisms fixing the diagonal idempotents.
\end{example}

\begin{example}\label{ex:hilbert-dual-An}
For $A=\Bbbk A_2$, Examples~\ref{ex:hilbert-triangular} and~\ref{ex:An-quadratic-dual} give
\[
h_A(t)=\begin{pmatrix}1&t\\0&1\end{pmatrix},\qquad
h_{A^!}(-t)^{\mathsf T}=\begin{pmatrix}1&-t\\0&1\end{pmatrix},\qquad
h_A(t)h_{A^!}(-t)^{\mathsf T}=I_2.
\]
The transpose aligns the indices because an $S$-dual reverses the two idempotents in \eqref{eq:dual-basis-idempotents}.
\end{example}
                   
\begin{example}\label{ex:alternating-quadratic-dual}
Proposition~\ref{prop:monomial-dual} exchanges the forbidden and surviving composable pairs.
Let $S=\Bbbk e_1\oplus\Bbbk e_2$ with coordinate idempotents, and let the bimodule $V=\Bbbk a\oplus\Bbbk b$ satisfy $e_1ae_2=a$, $e_2be_1=b$, with all other idempotent actions zero.
Take $A=T_S(V)/(ab)$.
The only composable words of length $2$ are $ab$ and $ba$. Thus
\[
A^!=T_S(D(V))/(a^*b^*),\qquad
A^!=S\oplus\Bbbk a^*\oplus\Bbbk b^*\oplus\Bbbk b^*a^*.
\]
The original relation $ab=0$ corresponds to the surviving word $b^*a^*$; the surviving word $ba$ corresponds to the dual relation $a^*b^*=0$.
In the order $e_1,e_2$, their Hilbert series are
\[
h_A(t)=\begin{pmatrix}1&t\\t&1+t^2\end{pmatrix},\qquad
h_{A^!}(t)=\begin{pmatrix}1+t^2&t\\t&1\end{pmatrix}.
\]
The degree $2$ terms occur in different entries: $ba\in e_2Ae_2$, whereas $b^*a^*\in e_1A^!e_1$.
\end{example}
\end{advanced}

\subsection{Symmetric and exterior algebras}

\begin{proposition}\label{prop:symmetric-exterior-dual}
For finite-dimensional $\Bbbk$-vector spaces $V,W$,
\[
\SymAlg(V)^!\cong\bigwedge D(V),\qquad
\bigl(\bigwedge W\bigr)^!\cong\SymAlg(D(W)).
\]
\end{proposition}

\begin{proof}
For $\SymAlg(V)$, choose a basis $x_1,\ldots,x_n$ of $V$, with dual basis $x_1^*,\ldots,x_n^*$ of $D(V)$.
Here $S=\Bbbk$, and \eqref{eq:D-tensor-duals} reads
\[
\theta(x_j^*\otimes x_i^*)(x_p\otimes x_q)=x_i^*(x_p)\,x_j^*(x_q)=\delta_{ip}\delta_{jq}.
\]
A tensor $\sum_{i,j=1}^n c_{ij}x_i^*\otimes x_j^*$ vanishes on each $x_p\otimes x_q-x_q\otimes x_p$ precisely when $c_{pq}=c_{qp}$ for $1\leq p<q\leq n$.
Thus the orthogonal complement of the symmetric relations is spanned by $(x_i^*)^2$ for $1\leq i\leq n$ and $x_i^*x_j^*+x_j^*x_i^*$ for $1\leq i<j\leq n$, which are the exterior relations on $D(V)$.

For $\bigwedge W$, choose a basis $y_1,\ldots,y_m$ of $W$, with dual basis $y_1^*,\ldots,y_m^*$ of $D(W)$.
A tensor $\sum_{i,j=1}^m c_{ij}y_i^*\otimes y_j^*$ annihilates the exterior relations precisely when $c_{ii}=0$ for $1\leq i\leq m$ and $c_{pq}+c_{qp}=0$ for $1\leq p<q\leq m$.
The orthogonal complement is therefore spanned by $y_i^*y_j^*-y_j^*y_i^*$ for $1\leq i<j\leq m$, which are the symmetric relations on $D(W)$.
The presentations in \eqref{eq:polynomial-exterior} give both isomorphisms.
\end{proof}

Consequently, to pair $\SymAlg(V)$ with $\bigwedge W$ by quadratic duality, take $W=D(V)$ and $y_i=x_i^*$, so $m=n$.
The $\Bbbk$-linear duality $D(D(V))\cong V$ from Example~\ref{ex:linear-duality} then identifies $(\bigwedge W)^!$ with $\SymAlg(V)$.

The proof works in characteristic $2$: the relations $(x_i^*)^2=0$ remain part of the exterior presentation even when the signs $+$ and $-$ coincide.
Example~\ref{ex:symmetric-orthogonals} is the calculation for two generators.
Quadratic duality need not preserve Hilbert series: if $\dim_\Bbbk V>0$, Example~\ref{ex:hilbert-polynomial-exterior} gives different series for $\SymAlg(V)$ and $\SymAlg(V)^!\cong\bigwedge D(V)$.

\begin{advanced}
The \emph{Koszul condition}, to be defined in Lecture~5, implies the identity
\begin{equation}\label{eq:matrix-koszul-preview}
h_A(t)h_{A^!}(-t)^{\mathsf T}=I_n
\qquad\text{when }A_0=\bigoplus_{i=1}^n\Bbbk e_i.
\end{equation}
This is \cite[Lemma~2.11.1]{BGS96}; the proof will be given in Lecture~8.
Equation~\eqref{eq:matrix-koszul-preview} is not asserted for arbitrary quadratic algebras.
For $\SymAlg(V)$ and its quadratic dual $\bigwedge D(V)$, with $n=\dim_\Bbbk V$, the case $S=\Bbbk$ gives $(1-t)^{-n}(1-t)^n=1$.
\end{advanced}

\begin{advanced}[suppressfirstparskip=true,innertopmargin=0pt]
\subsection{Tensor products and signs}

\begin{definition}\label{def:algebra-tensor-products}
For graded $\Bbbk$-algebras $A,B$, the \defn{ordinary tensor product} $A\otimes_\Bbbk B$ and the \defn{graded tensor product} $A\otimes_\Bbbk^{\mathrm{gr}}B$ have underlying vector space $A\otimes_\Bbbk B$ and the total grading of Definition~\ref{def:total-tensor-grading}, with products
\begin{equation}\label{eq:signed-tensor-product}
\begin{aligned}
(a\otimes b)(a'\otimes b')&=aa'\otimes bb'
&&\text{in }A\otimes_\Bbbk B,\\
(a\otimes b)(a'\otimes b')&=(-1)^{\deg b\,\deg a'}aa'\otimes bb'
&&\text{in }A\otimes_\Bbbk^{\mathrm{gr}}B
\end{aligned}
\end{equation}
for homogeneous $a,a'\in A$ and $b,b'\in B$, extended bilinearly.
\end{definition}

Both products have identity $1_A\otimes1_B$.
For the graded product, associativity follows from
$\deg b\,\deg a'+(\deg b+\deg b')\deg a''=\deg b'\,\deg a''+\deg b(\deg a'+\deg a'')$
for homogeneous $a',a''\in A$ and $b,b'\in B$.
For homogeneous $a\in A$ and $b\in B$, the elements $a\otimes1$ and $1\otimes b$ commute in the ordinary tensor product; in the graded tensor product they satisfy
\[
(1\otimes b)(a\otimes1)=(-1)^{\deg a\,\deg b}(a\otimes1)(1\otimes b).
\]
Suppose $A_0=S$ and $B_0=T$ are finite-dimensional semisimple and $S\otimes_\Bbbk T$ is semisimple.
The sign does not change the underlying graded $(S\otimes_\Bbbk T)$-bimodule, so the two tensor products have the same Hilbert series whenever $A,B$ are locally finite-dimensional.
For $S=\Bbbk^{\oplus n}$ and $T=\Bbbk^{\oplus m}$, Lemma~\ref{lem:hilbert-field-tensor} gives its entries:
\[
h_{A\otimes_\Bbbk B}(t)_{(i,p),(j,q)}
=h_A(t)_{ij}h_B(t)_{pq}.
\]
In particular, when $S=T=\Bbbk$, this is $h_A(t)h_B(t)$.

\begin{proposition}\label{prop:dual-tensor-products}
Let $A=T_S(V)/(R)$ and $B=T_T(W)/(Q)$ be quadratic algebras, and assume that $S\otimes_\Bbbk T$ is semisimple.
Then both tensor products are quadratic over $S\otimes_\Bbbk T$, and
\begin{equation}\label{eq:dual-tensor-products}
(A\otimes_\Bbbk B)^!\cong A^!\otimes_\Bbbk^{\mathrm{gr}}B^!,\qquad
(A\otimes_\Bbbk^{\mathrm{gr}}B)^!\cong A^!\otimes_\Bbbk B^!.
\end{equation}
\end{proposition}

The mixed commuting relations become anticommuting relations under the reversed tensor pairing.

The assumption on $S\otimes_\Bbbk T$ holds when $S,T$ are products of matrix algebras over $\Bbbk$, since $\Mat_a(\Bbbk)\otimes_\Bbbk\Mat_b(\Bbbk)\cong\Mat_{ab}(\Bbbk)$.
It does not follow from semisimplicity alone: take $S=T$ to be the field in Example~\ref{ex:inseparable-envelope}.

\begin{proof}
The degree $1$ part of $A\otimes_\Bbbk B$ is $(V\otimes_\Bbbk T)\oplus(S\otimes_\Bbbk W)$.
By Lemma~\ref{lem:tensor-properties}, its tensor square over $S\otimes_\Bbbk T$ decomposes into the four summands
\[
(V\otimes_S V)\otimes_\Bbbk T
\ \oplus\ \underbrace{V\otimes_\Bbbk W}_{\text{order }VW}
\ \oplus\ \underbrace{V\otimes_\Bbbk W}_{\text{order }WV}
\ \oplus\ S\otimes_\Bbbk(W\otimes_T W).
\]
Its quadratic relations consist of $R\otimes_\Bbbk T$, $S\otimes_\Bbbk Q$, and the mixed relations
\begin{equation}\label{eq:mixed-tensor-relations}
(v\otimes1)(1\otimes w)-(1\otimes w)(v\otimes1)
\quad(v\in V,\ w\in W).
\end{equation}
Here the products are taken in the tensor algebra on the degree $1$ part, and their span is an $(S\otimes_\Bbbk T)$-sub-bimodule.
Let $C$ be the quotient of that tensor algebra by the three specified relation spaces.
The generator maps induce algebra homomorphisms $A\to C$, $a\mapsto\overline a$, and $B\to C$, $b\mapsto\overline b$.
Their images commute: the degree $0$ actions already commute across the two factors, and \eqref{eq:mixed-tensor-relations} makes the degree $1$ generators commute.
Consequently $A\otimes_\Bbbk B\to C$, $a\otimes b\mapsto\overline a\,\overline b$, is an algebra homomorphism.
It is inverse to the canonical homomorphism $C\to A\otimes_\Bbbk B$, since both composites fix the generating copies of $S,T,V,W$. Thus the stated relations give exactly $A\otimes_\Bbbk B$.
For the graded tensor product, the same construction uses $\overline b\,\overline a=(-1)^{\deg a\deg b}\overline a\,\overline b$ for homogeneous $a,b$.

We use the description of quadratic duals by right $S$-duals (advanced material in Section~\ref{sec:quadratic-dual}), which applies to the semisimple algebra $S\otimes_\Bbbk T$ and agrees with Definition~\ref{def:quadratic-duals} when $S\otimes_\Bbbk T$ is a product of copies of $\Bbbk$.
The right $(S\otimes_\Bbbk T)$-duals of the two generator summands are $V^*\otimes_\Bbbk T$ and $S\otimes_\Bbbk W^*$, with evaluations
\[
(f\otimes t)(v\otimes t')=f(v)\otimes tt',\qquad
(s\otimes g)(s'\otimes w)=ss'\otimes g(w).
\]
These identifications hold for finite free right modules and hence for their direct summands by Corollary~\ref{cor:semisimple-summands}.
The reversed tensor evaluation pairs each pure summand with its corresponding dual summand and exchanges the two mixed summands; all other pairings are zero.
The pure relation spaces have orthogonal complements $R^\perp\otimes_\Bbbk T$ and $S\otimes_\Bbbk Q^\perp$ in the corresponding pure summands.
The mixed relation space is the set of pairs $(z,-z)$ in the two copies of $V\otimes_\Bbbk W$.
For mixed tensors, the two reversed evaluations of $f\in V^*$ and $g\in W^*$ both give $f(v)\otimes g(w)$ on the corresponding copy.
Thus an annihilating pair of functionals must have equal components. As these simple tensors span the full dual, the mixed annihilator is spanned by
\[
(f\otimes1)(1\otimes g)+(1\otimes g)(f\otimes1)
\quad(f\in V^*,\ g\in W^*).
\]
These are the mixed relations of $A^!\otimes_\Bbbk^{\mathrm{gr}}B^!$, proving the first isomorphism in \eqref{eq:dual-tensor-products}.
For $A\otimes_\Bbbk^{\mathrm{gr}}B$, the minus sign in \eqref{eq:mixed-tensor-relations} is a plus sign. The mixed annihilator then has a minus sign, giving the ordinary tensor product on the dual side.
\end{proof}

\begin{example}\label{ex:commuting-square-dual}
We compute how Proposition~\ref{prop:dual-tensor-products} changes a commutation relation in a matrix example.
Take $A=\Bbbk A_2\otimes_\Bbbk\Bbbk A_2$.
Its degree $0$ algebra has basis $e_{ii}\otimes e_{jj}$ for $i,j\in\{1,2\}$, and its degree $1$ part has basis
\[
a=e_{12}\otimes e_{11},\quad
b=e_{22}\otimes e_{12},\quad
c=e_{11}\otimes e_{12},\quad
d=e_{12}\otimes e_{22}.
\]
By Lemma~\ref{lem:idempotent-tensors}, the tensor square of the degree $1$ part has basis $ab,cd$, and both map to $e_{12}\otimes e_{12}$. Thus the relation space is $\Bbbk(ab-cd)$.
Using the reversed generators from \eqref{eq:dual-basis-idempotents}, its orthogonal complement is $\Bbbk(b^*a^*+d^*c^*)$.
Hence the dual presentation has an anticommuting relation:
\[
ab=cd\text{ in }A,\qquad
b^*a^*=-d^*c^*\text{ in }A^!.
\]
Order the degree $0$ idempotents as
$e_{11}\otimes e_{11},e_{22}\otimes e_{11},e_{11}\otimes e_{22},e_{22}\otimes e_{22}$.
The homogeneous bases give
\[
h_A(t)=\begin{pmatrix}
1&t&t&t^2\\
0&1&0&t\\
0&0&1&t\\
0&0&0&1
\end{pmatrix},\qquad
h_{A^!}(t)=h_A(t)^{\mathsf T}.
\]
The signs agree with Proposition~\ref{prop:dual-tensor-products}; in characteristic $2$, the same presentations remain valid with $-1=1$.
Here replacing $d^*$ by $-d^*$ removes the minus sign. In this example the sign distinguishes the specified presentations, not the isomorphism classes of algebras.
\end{example}

\begin{example}\label{ex:essential-tensor-sign}
The symmetric/exterior pair shows that the sign in a graded tensor product cannot always be omitted.
Let $V=\Bbbk x_1\oplus\Bbbk x_2$ and let $y_1,y_2$ be the dual basis of $W=D(V)$. Proposition~\ref{prop:dual-tensor-products} gives
\[
\SymAlg(V)^!
\cong\bigl(\Bbbk[x_1]\otimes_\Bbbk\Bbbk[x_2]\bigr)^!
\cong\frac{\Bbbk[y_1]}{(y_1^2)}\otimes_\Bbbk^{\mathrm{gr}}
       \frac{\Bbbk[y_2]}{(y_2^2)}
\cong\bigwedge W.
\]
If $\Char\Bbbk\neq2$, the ordinary tensor product of the two square-zero algebras is commutative, whereas the graded tensor product is not: $y_1y_2=-y_2y_1\neq0$.
They are therefore not isomorphic. This sign cannot be removed by a change of generators.
\end{example}
\end{advanced}

\subsection{Exercises}

\begin{exercise}
For $A=\Bbbk A_4/(e_{13})$, write a presentation of $A^!$ using reversed generators and find homogeneous bases of $A$ and $A^!$.
Check that $(A^!)^!\cong A$ by computing the relations of $A^!$ directly.
\emph{Hint.} Among $e_{12}e_{23}$ and $e_{23}e_{34}$, only the first is a relation. Complement the forbidden pairs and reverse their order.
\end{exercise}

\begin{exercise}
Let $S=\bigoplus_{i=1}^n\Bbbk e_i$ with coordinate idempotents, and let ${}_SV_S$ be a finite-dimensional bimodule in degree $1$.
Set $C_{ij}=[e_iV:e_jS]=\dim_\Bbbk(e_iVe_j)$ for $1\leq i,j\leq n$, and put
\[
A=T_S(V)/(V\otimes_S V)=S\oplus V.
\]
\begin{enumerate}
\item Prove that
\[
A^!=T_S(D(V)),\qquad h_A(t)=I_n+tC,\qquad
h_{A^!}(t)=(I_n-tC^{\mathsf T})^{-1}.
\]
\item Deduce that $h_A(t)h_{A^!}(-t)^{\mathsf T}=I_n$.
\end{enumerate}
\emph{Hint for (1).} Find $(V\otimes_S V)^\perp$ and compute the multiplicities $[e_iD(V):e_jS]$.
\end{exercise}

\begin{advanced}
\begin{exercise}
Let $S=\bigoplus_{i=1}^4\Bbbk e_i$ with coordinate idempotents, and let ${}_SV_S$ have basis $\{a,b,c,d\}$ with
\[
a\in e_1Ve_2,\qquad b\in e_2Ve_4,\qquad
c\in e_1Ve_3,\qquad d\in e_3Ve_4,
\]
and $e_iVe_j=0$ for all other pairs.
For $\lambda\in\Bbbk$, put $R_\lambda=\Bbbk(ab-\lambda cd)\subset V\otimes_S V$ and $A_\lambda=T_S(V)/(R_\lambda)$.
\begin{enumerate}
\item Using $\theta$, compute $R_\lambda^\perp$ and a presentation of $A_\lambda^!$ for every $\lambda$, including $\lambda=0$.
\item Compute $(A_\lambda^!)^!$ and recover the relation $ab-\lambda cd$ under the canonical identification $D(D(V))\cong V$.
\item Prove that $A_\lambda\cong A_\mu$ as graded algebras if and only if $\lambda,\mu$ are both zero or both non-zero.
\end{enumerate}

\emph{Hint for (1).} Evaluate $\alpha b^*a^*+\beta d^*c^*$ on $ab-\lambda cd$ for $\alpha,\beta\in\Bbbk$.

\emph{Hint for (3).} A graded isomorphism permutes the coordinate idempotents of $S$.
The degree $1$ components force it to fix $e_1,e_4$ and either fix or interchange $e_2,e_3$.
Compare which of the multiplication maps
\[
e_1(A_\lambda)_1e_i\otimes_\Bbbk e_i(A_\lambda)_1e_4
\longrightarrow e_1(A_\lambda)_2e_4
\quad\text{for }i\in\{2,3\}
\]
are zero. For non-zero $\lambda,\mu$, rescale $d$.
\end{exercise}
\end{advanced}

